\section{Open questions}

The following five probes were flagged during replication as concrete
next-step experiments that this study does \emph{not} settle. Each has a
basis in what was (and was not) exercised, and a specific CPU-scale
follow-up that stays within free-endpoint budgets.

\begin{enumerate}

\item \textbf{Does RQAE's quadratic scaling advantage survive under a
realistic gate/readout noise model at the depths implied by $k_{\max}$
for $\varepsilon = 10^{-3}$, or does noise flatten the slope back toward 2?}
\emph{Basis:} this replication used a noiseless \texttt{AerSimulator}
with real shot Bernoulli sampling only; noise robustness was not tested.
\emph{Next steps:} extend \texttt{code/rqae.py} to accept an
\texttt{AerSimulator} with \texttt{NoiseModel.from\_backend(FakeSherbrooke)}
or a parameterized depolarizing + readout model; re-run the 12-config sweep
at 3--5 noise-strength levels; refit the log--log slope $\alpha$ and check
whether $\alpha$ crosses back above $1.5$ as noise increases.

\item \textbf{Is RQAE's constant on $N_{\text{oracle}}$ competitive with
MLQAE (Suzuki 2019) and IQAE (Grinko 2019) at matched
$(\varepsilon, \gamma)$, or is the only genuine advantage of RQAE the
sign-recovery capability?} \emph{Basis:} we benchmarked RQAE only against
the classical Hoeffding baseline, not against other QAE variants.
\emph{Next steps:} implement MLQAE and IQAE on the same single-qubit toy
oracle over the same $(\varepsilon, \gamma, a_{\text{true}})$ grid; report
mean $N_{\text{oracle}}$ for each method at $\varepsilon \in \{0.05, 0.01, 0.005\}$
using \texttt{qiskit-algorithms} reference implementations.

\item \textbf{Does the RQAE framework extend cleanly to multi-amplitude /
vector estimation, or does the shift-and-amplify trick fundamentally
require a single scalar amplitude?} \emph{Basis:} the paper and this
replication address a single real amplitude; modern QAE applications often
need $k$ amplitudes simultaneously. \emph{Next steps:} prototype a
2-amplitude 2-qubit oracle; run RQAE independently on each and jointly and
compare total $N_{\text{oracle}}$ to $k \cdot N_{\text{oracle,single}}$.

\item \textbf{On a specific published option-pricing benchmark (e.g.\
European call with piecewise-linear payoff, Woerner--Egger 2019 setup),
how does RQAE compare in oracle-call count to canonical QAE at matched
target precision on the expected payoff?} \emph{Basis:} the replication
used a single-qubit toy oracle, not a realistic multi-qubit finance
$A$-operator. \emph{Next steps:} wire a \texttt{qiskit-finance}
\texttt{EuropeanCallPricingObjective} into \texttt{rqae.py} as the $A$
operator; re-run at $\varepsilon \in \{10^{-2}, 5 \times 10^{-3}\}$ for
10 reps and compare against the QAE-based pricer in the same repo.

\item \textbf{Does an adaptive-schedule RQAE variant (choosing $k_{i+1}$
from the observed $\hat{p}_i$ posterior rather than the worst-case
Eq.~(14)) achieve empirically better constants without breaking coverage?}
\emph{Basis:} Eq.~(14) is worst-case; the observed $\hat{p}_i$ is often
sharper than the Hoeffding interval, but the paper does not explore an
adaptive rule. \emph{Next steps:} add an \texttt{adaptive\_rqae()} variant
using a Beta-Binomial / $\sin(2k+1)\theta$ posterior for $k_{i+1}$; re-run
the same 12-config sweep at 25 reps with matched seeds for a paired
comparison of mean $N_{\text{oracle}}$ and coverage.

\end{enumerate}
